\section{Vorhersage kündigungsgefährdeter KV-Kundinnen und -Kunden}
\label{sec:churn}

\subsection{Zielsetzung und fachlicher Kontext}

Der zweite Use Case unterstützt Vertrieb und Kundenbetreuung in der privaten
Krankenversicherung. Ein Modell soll die Kündigungswahrscheinlichkeit von
Bestandskundinnen und -kunden bestimmen, wesentliche Einflussfaktoren
ausweisen und gefährdete Fälle für Rückhaltemaßnahmen priorisieren. Als
fachliches Ziel nennt das Anforderungsszenario eine Verringerung der
Kündigungsquote um 10\,\% gegenüber dem Vorjahreszeitraum
\citep{spiteri2018,nordakademieSzenarioVersicherung}. Dieser Zielwert ist im
späteren Betrieb zu validieren und stellt kein bereits erzieltes Ergebnis dar.

\subsection{Architekturrelevante Anforderungen}

Die Prognose nutzt vor allem Vertragsdaten zu Laufzeit, Beitragsentwicklung
und Tarifwechseln; historische Kündigungen und Neuabschlüsse liefern die
Trainingsgrundlage. Kunden-, CRM- und Zahlungsdaten bilden bisherige
Interaktionen, Beschwerden und Zahlungsverhalten ab. Ergänzend werden
aggregierte Leistungsabrechnungen sowie sozioökonomische Kontextdaten
einbezogen.

Das vollständige Datenprofil mit Struktur, Aktualität, Speicherdauer und
Analysebedarf ist in Tabelle~\ref{tab:daten-churn} im Anhang dargestellt
\citep[Folien~17--18]{nordakademieSzenarioVersicherung}.

Die Quellen sind überwiegend strukturiert und werden kontinuierlich oder
hochaktuell aktualisiert. Für Training und Ursachenanalyse müssen historische
Stände erhalten bleiben; die Kundenbetreuung benötigt aktuelle
Modellergebnisse. Gesundheitsbezogene Leistungsdaten sind zu minimieren und
streng zu schützen. Training und Anwendung des Modells müssen innerhalb des
jeweils zulässigen regionalen Datenraums erfolgen
\citep{euGDPR}.

Aus den fachlichen Zielen und den unterschiedlichen Anforderungen an
Struktur, Aktualität, Schutz und Historisierung der Daten werden in
Kapitel~\ref{chap:architektur} die benötigten Architekturkomponenten
abgeleitet und zu einer ganzheitlichen Zielarchitektur zusammengeführt.
